\documentclass[11pt]{article}

% Shared preamble for the SKJS soundness corpus (main paper + deep dives).
% Loaded via \input in each document's preamble. Documents may define
% additional packages, theorem environments, and macros AFTER this input,
% but must not redefine anything here.
\usepackage{amsmath}
\usepackage{amssymb}
\usepackage{amsthm}
\usepackage{mathtools}
\usepackage{stmaryrd}
\usepackage[margin=1in]{geometry}
\usepackage{tcolorbox}
\usepackage{booktabs}
\usepackage{array}
\usepackage{longtable}
\usepackage[hyphens]{url}
\urlstyle{tt}

% Long typewriter tokens (paths, identifiers) defeat hyphenation;
% allow inter-word stretch instead of overfull lines.
\emergencystretch=2em

\newtheorem{theorem}{Theorem}
\newtheorem{lemma}[theorem]{Lemma}
\newtheorem{definition}[theorem]{Definition}
\newtheorem{corollary}[theorem]{Corollary}
\newtheorem{axiom}[theorem]{Axiom}
\newtheorem{obligation}[theorem]{Obligation}

% Summary box for key definitions
\newtcolorbox{keydef}{
  colback=gray!5,
  colframe=gray!50,
  boxrule=0.5pt,
  arc=2pt,
  left=4pt,right=4pt,top=2pt,bottom=2pt
}

% Core semantic notation shared by every document
\newcommand{\sem}[1]{\llbracket #1 \rrbracket}
\newcommand{\Val}{\mathsf{Val}}

% Implementation anchor: file / class / function at the pinned revision
\newcommand{\impl}[1]{\par\noindent{\footnotesize\textsf{impl:} \texttt{#1}}}

% Single source for the modeled revision. Every document's \date uses this;
% the ledger and tooling read the same values.
\newcommand{\skjspinhash}{9108f339}
\newcommand{\skjspindate}{2026-07-02}
\newcommand{\skjspin}{Models revision \texttt{\skjspinhash} of \texttt{main} (\skjspindate)}


\usepackage[T1]{fontenc}

\emergencystretch=1.5em
\tolerance=2000

\newtheorem{finding}[theorem]{Finding}

% Macros copied from the main paper (skjs-type-soundness.tex)
\newcommand{\Pred}{\mathcal{P}(\Val)}
\newcommand{\TType}{\mathsf{TType}}
\newcommand{\Clauses}{\mathsf{Clauses}}
\newcommand{\STypes}{\mathsf{STypes}}
\newcommand{\Subst}{\mathsf{Subst}}
\newcommand{\SVOpen}{\mathsf{SVOpen}}
\newcommand{\SVClosed}{\mathsf{SVClosed}}
\newcommand{\Env}{\Gamma}

% Local macros
\newcommand{\env}[2]{\mathcal{E}_{#1}(#2)}   % envelope of a variable
\newcommand{\stle}{\sqsubseteq}              % pointwise state order

\title{The Substitution State Machine: Monotonicity, Monotone Reads,
Clone Stability, and Rollback\\
{\large Deep dive E for ``Refinement and Inference Growth in SKJS''}}
\author{}
\date{\skjspin}

\begin{document}

\maketitle

\begin{abstract}
The main paper models the inference substitution as a state machine;
four of its ledger rows concern that machine's state discipline —
axiom \textbf{L-MONO} (refinement-time monotonicity), the
obligation \emph{Monotone Reads Under Growth}, assumption
\textbf{L-CLONE} (clone stability), and the obligation
\emph{No Commitment on Failure} — and this paper is the evidence behind
their current table statuses. This deep dive completes
the state model the main paper elides — the version counter, the four
caches, the environment-side cycle sets, constraint attachment, and the
candidate-freezing flag — and works the four rows as one coherent
discipline. The load-bearing discovery is representational:
\texttt{TVar::setConstraint} is a \emph{functional update} on a frozen
value, not a mutation of shared state, which removes the envelope-shrinking
threat the main paper flagged for \textbf{L-MONO} and \textbf{L-CLONE}.
With that, \textbf{L-MONO} is proved under two stated modeling conditions;
\emph{Monotone Reads} is proved by a read-variance audit everywhere except
one \textsc{violated} residue, MR-1 — a freeze-window unsealed read
memoized in \texttt{seenSub}, witnessed end-to-end by an order-dependent
unsound acceptance (issue draft
\path{issues/freeze-window-stale-memo-caches.md}); \textbf{L-CLONE} is
proved conditional on a strengthened, fork-state form of \textbf{L-INTER}
— the clone-generated-identifier hazard is verified \emph{unreachable}
(no \texttt{gen*} call site lies in the clone-compute-install tree),
guarded by a recommended CI assertion; and \emph{No Commitment on
Failure} is discharged by an exhaustive save/restore audit with one
intentional, separately-argued exception (array-literal form commitments
persisting on failed checks). Failed proof attempts are reported as
findings.
\end{abstract}

\tableofcontents

%==========================================================================
\section{Purpose and Rows Addressed}
\label{sec:purpose}

This document supports the ``Assumptions and Proof Obligations'' ledger of
the main paper. It addresses four rows --- \textbf{L-MONO}, Monotone Reads
Under Growth, \textbf{L-CLONE}, and No Commitment on Failure --- treated in
Sections~\ref{sec:lmono}, \ref{sec:monoreads}, \ref{sec:lclone},
and~\ref{sec:rollback}, and summarized row by row in \S\ref{sec:status}.

All file and line references are against commit \texttt{\skjspinhash}. Line
numbers are cited as \texttt{file:line}; they identify the revision
precisely and will drift with edits.

%==========================================================================
\section{The Complete State Model}
\label{sec:model}

The main paper's Definition ``Substitution States'' records only the
binding map and the freeze flag. The implementation state is larger. This
section is the reference model; everything later in the paper is proved
against it.

\subsection{Saved state: \texttt{SubstState}}

\begin{definition}[Full Substitution State]
\label{def:fullstate}
A substitution state $\rho$ consists of the fields of the Skip class
\texttt{SubstState} (\texttt{TTypeSubst.sk:13--49}):
\begin{itemize}
\item $\mathsf{values} : \mathsf{TVarID} \rightharpoonup
  \SVOpen(\TType) \mid \SVClosed(\TType)$ — the binding map of the main
  paper's transition table;
\item $\mathsf{version} \in \mathbb{N}$ — a counter with two duties:
  (a) it is incremented immediately before \emph{every} write to
  $\mathsf{values}$ (\texttt{grow}, \texttt{TTypeSubst.sk:181};
  \texttt{get}'s two sealing writes, \texttt{:195} and \texttt{:205}); and
  (b) it doubles as the fresh-identifier supply — \texttt{gen},
  \texttt{genForField}, and \texttt{genForLiteralCommitment}
  (\texttt{:237--253}) each return \texttt{TVarID(version)} and increment.
  Consequence: within one linear execution segment, distinct
  $\mathsf{values}$ states have distinct versions
  (Lemma~\ref{lem:version-values});
\item $\mathsf{origins}$ — diagnostic provenance for variables
  (\texttt{setOrigin}); no semantic content;
\item four caches, all part of the saved state:
  \begin{itemize}
  \item $\mathsf{seenSub} \subseteq \TType \times \TType$ — a
    \emph{success memo} for \texttt{TType::isSub}: a pair is inserted on
    acceptance (\texttt{TTypeSubtyping.sk:1591}) and consulted before
    re-checking (\texttt{:1588}). Not version-keyed;
  \item $\mathsf{seenInter}, \mathsf{seenUnion} :
    (\TType \times \TType) \rightharpoonup \TType$ — result memos for
    \texttt{TType::inter} (\texttt{TTypeInter.sk:107--125}; pairs
    containing a bare variable are excluded, \texttt{:105}) and
    \texttt{TType::union} (\texttt{TTypeUnion.sk:29--32});
  \item $\mathsf{seenExpansion} : \mathsf{TSInst} \rightharpoonup \TType$
    — memo of one-level named-type expansion
    (\texttt{TType.sk:650--653}); the cached body is \emph{symbolic} in
    the type arguments (arguments occur as free variables, re-resolved at
    every use), which is why reuse under growth is sound
    (comment \texttt{TTypeSubst.sk:17--30}).
  \end{itemize}
\end{itemize}
\texttt{save()} snapshots all seven fields into a frozen
\texttt{SavedSubst}; \texttt{restore()} reinstates all seven
(\texttt{TTypeSubst.sk:36--48}). Because every field is an immutable
persistent map/set, \texttt{save} is $O(1)$ and lossless.
\impl{TTypeSubst.sk / SubstState, SavedSubst}
\end{definition}

\subsection{Unsaved state}

Three pieces of mutable or threaded state live \emph{outside}
\texttt{SubstState}; each needs its own safety argument.

\paragraph{(U1) \texttt{freezeCandidates}.}
Declared on \texttt{Subst} itself, not on \texttt{SubstState}
(\texttt{TTypeSubst.sk:76--89}), with an explicit comment that
save/restore never touch it; it is push/popped around the
constraint-admissibility check in \texttt{growTVars}
(\texttt{TTypeSubtyping.sk:1431--1443}). Verified by inspection: those two
lines are the \emph{only} writes to the flag in \texttt{src/}, the
push/pop brackets a straight-line region with no early return (the
\texttt{break} at \texttt{:1440} exits only the inner constraint loop),
and nested checks restore the outer value via \texttt{wasFreezing}. So
excluding it from \texttt{SavedSubst} is safe: a restore can never observe
a torn flag, because every save/restore pair either lies entirely inside
or entirely outside a freeze window.

\paragraph{(U2) \texttt{tenv.seenSub} and \texttt{tenv.seenInter}.}
\texttt{TEnv} is a \emph{frozen} Skip class (\texttt{TEnv.sk:12}); the
apparent mutations \texttt{!tenv.seenSub = \ldots}
(\texttt{TTypeSubtyping.sk:1537}) and \texttt{!tenv.seenInter = \ldots}
(\texttt{TTypeInter.sk:157}) are functional updates of the \emph{local}
binding, propagated only into callees that receive the updated
\texttt{tenv}. These sets are therefore \emph{path-scoped coinduction
hypotheses}, not caches: an entry is visible exactly in the dynamic
extent of the recording activation and vanishes when it returns. This
structural fact does most of the work in the \#446 argument
(\S\ref{sec:tenv-seensub}). \texttt{tenv.seenSub} is keyed by
$(\mathsf{version}, u_1, u_2)$ (\texttt{TEnv.sk:74});
\texttt{tenv.seenInter} by the bare pair (\texttt{TEnv.sk:69}), which is
safe because its only effect is to return an \emph{unexpanded}
intersection on a cycle (\texttt{TTypeInter.sk:154--156}) — a
representation choice, not a judgement.

\paragraph{(U4) \texttt{tenv.expandInter}.}
A mode flag on \texttt{TEnv} selecting full versus deferred
(\texttt{unexpandedInter}) field intersection; set by
\texttt{TInter::expand} and, like \texttt{tenv.seenSub}, threaded through
the frozen-\texttt{TEnv} update discipline, so the setting should be
path-scoped. This flag was found missing from this inventory by external
review; its path-scoping and the containment coverage of the deferred
form are the \textsc{open} registry row TENV-EXPANDINTER.

\paragraph{(U3) \texttt{results} and \texttt{context}.}
The diagnostics/callgraph accumulator \texttt{results} is shared by
\texttt{clone()} (\texttt{TTypeSubst.sk:137--149}) and never rolled back
(e.g.\ \texttt{recordSubtypingEdge} on \texttt{TFun::isSub} success,
\texttt{TTypeSubtyping.sk:1681--1688}). It carries no typing state, so it
is outside $\rho$; \S\ref{sec:lclone} flags the one place this is worth a
remark.

\subsection{Constraint attachment is a value operation}
\label{sec:setconstraint}

A natural worry about the main paper's Axiom \textbf{L-MONO} would be
that adding a constraint to a \emph{shared} variable
(\texttt{TVar::setConstraint}) shrinks that variable's envelope, so an
arbitrary type mentioning it could lose members. The premise of that
worry is false — establishing this is the single most consequential fact
in this paper, and the axiom's remark in the main paper records it.

\begin{lemma}[setConstraint is Functional]
\label{lem:setconstraint}
\texttt{TVar::setConstraint(ty)} does not modify any state observable by
another type term. It returns a \emph{new} \texttt{TVar} value carrying
the enlarged constraint set; every pre-existing occurrence of the same
variable identifier retains its own constraint set.
\end{lemma}
\begin{proof}[Proof (by inspection of the pinned source)]
\texttt{TVar} is a frozen (non-\texttt{mutable}) class and
\texttt{setConstraint} a non-\texttt{mutable} method
(\texttt{TType.sk:2589} and \texttt{:2590--2591}):
\begin{center}
\texttt{fun setConstraint(ty: TType): TVar \{}\\
\texttt{\ \ !this.constraints = this.constraints.set(ty); this \}}
\end{center}
In Skip, field assignment through \texttt{!this.f = e} in a frozen-class
method is a \emph{functional record update of the local binding}
\texttt{this}; the method must return the updated copy for the change to
be visible at all. The same idiom appears throughout the pinned revision
on unambiguously frozen classes
(\texttt{STErrorDetail::ensureIntermediateDepth},
\texttt{TTypeSubtyping.sk:40--45}; \texttt{TEnv::instantiate},
\texttt{TEnv.sk:131}). Two corroborations: (i) \texttt{constraints}
participates in \texttt{TVar::xxhash} (\texttt{TType.sk:2598--2607}) and
in \texttt{Orderable} equality — a shared-mutation semantics would break
hashing of interned type terms; (ii) both call sites use the return value
(\texttt{TType::var(tvar.setConstraint(ty2))},
\texttt{TTypeInter.sk:145--151}; \texttt{!tvar = tvar.setConstraint(...)},
\texttt{TEnv.sk:314}) — a genuine mutation would make that redundant.
\end{proof}

\begin{corollary}[Occurrence-Relative Envelopes]
\label{cor:occurrence}
The ``constraint intersection'' step of \texttt{Subst::get}
(\texttt{TTypeSubst.sk:218--221}) intersects the binding with the
constraints \emph{of the occurrence being read}. A variable identifier
therefore has no global constraint set; constraints are part of the type
term. Refinement, intersection, or any other operation attaching a
constraint produces a new term and cannot shrink the denotation of any
type already installed in an environment. The state $\rho$ of
Definition~\ref{def:fullstate} is the \emph{complete} shared mutable
typing state.
\end{corollary}

\subsection{The extended transition table}

\begin{definition}[Extended Transitions]
\label{def:ext-transitions}
The operations that evolve $\rho$, extending the main paper's
Definition~6 (its rows are unchanged and not repeated):
\begin{center}
\renewcommand{\arraystretch}{1.2}
\footnotesize
\begin{tabular}{@{}lll@{}}
\toprule
op & effect on $\rho$ & anchor \\
\midrule
$\mathsf{grow},\ \mathsf{get}(\mathsf{cov})$ writes &
  per main paper; $\mathsf{version}{+}{+}$ before each write &
  \texttt{TTypeSubst.sk:151--222} \\
$\mathsf{gen}/\mathsf{genFor*}$ &
  fresh id $=$ $\mathsf{version}$; $\mathsf{version}{+}{+}$; no binding &
  \texttt{TTypeSubst.sk:237--253} \\
cache insert & add to one of the four memo maps; no version bump &
  \S\ref{def:fullstate} \\
$\mathsf{setOrigin}$ & origins only & \texttt{TTypeSubst.sk:271} \\
$\mathsf{save}/\mathsf{restore}$ & snapshot / reinstate all seven fields &
  \texttt{TTypeSubst.sk:36--48} \\
$\mathsf{clone}$ & fork: new \texttt{Subst}, same seven fields, shared
  \texttt{results} & \texttt{TTypeSubst.sk:137--149} \\
freeze push/pop & toggle \texttt{freezeCandidates} (unsaved) &
  \texttt{TTypeSubtyping.sk:1431--1443} \\
\bottomrule
\end{tabular}
\end{center}
Verified write-site completeness: the only assignments to
\texttt{values} in \texttt{src/} are in \texttt{grow}, \texttt{get}, and
\texttt{restore}; the only version writes are those plus \texttt{gen*};
external \texttt{grow} callers (\texttt{TArrayLitTyping.sk:147,352};
\texttt{TVariadicTyping.sk:125}; \texttt{TypeCheckUtils.sk:310,321,332})
go through the same transitions.
\end{definition}

Note that cache inserts do \emph{not} bump the version. This is sound for
the three result memos by their own stability arguments (sealing for
\texttt{seenInter}, \texttt{TTypeInter.sk:90--96}; symbolic arguments for
\texttt{seenExpansion}) but it makes the version-keying of
\texttt{tenv.seenSub} track \emph{binding} state only — exactly what the
coinduction hypothesis depends on (\S\ref{sec:tenv-seensub}).

%==========================================================================
\section{The Order-Theoretic Core}
\label{sec:order}

All four ledger rows reduce to statements about one partial order on
states and two lemmas about it.

\begin{definition}[Envelope; State Order]
\label{def:order}
For a state $\rho$ and identifier $\alpha$, the \emph{envelope} is
\[
\env{\rho}{\alpha} =
\begin{cases}
\emptyset & \alpha \notin \mathrm{dom}(\rho.\mathsf{values})\\
\sem{t}_\rho & \rho.\mathsf{values}(\alpha) \in
  \{\SVOpen(t), \SVClosed(t)\}.
\end{cases}
\]
Write $\rho \stle \rho'$ (``$\rho'$ extends $\rho$'') iff
$\env{\rho}{\alpha} \subseteq \env{\rho'}{\alpha}$ for every $\alpha$.
The choice $\emptyset$ for unbound variables is the modeling counterpart
of $\mathsf{get}(\mathsf{contra})$ returning $\bot$ without writing
(\texttt{TTypeSubst.sk:211--212}).
\end{definition}

\begin{definition}[Modeling Conditions]
\label{def:conditions}
\begin{itemize}
\item[\textbf{(M1)}] \emph{Semantic monotonicity.} For every type term
  $u$, $\sem{u}_\rho$ is monotone in the envelopes of the identifiers
  occurring in $u$: $\rho \stle \rho'$ implies
  $\sem{u}_\rho \subseteq \sem{u}_{\rho'}$. This holds for the main
  paper's semantics by structural induction, because every
  state-sensitive position is read positively: variable occurrences
  denote $\env{\rho}{\alpha} \cap \bigcap \sem{\mathrm{cstr}}$
  (Corollary~\ref{cor:occurrence}); field membership constrains only the
  \emph{get} type (main paper, object semantics, which makes set
  types opaque); call/new signatures are opaque to $\sem{\cdot}$ except
  for the callability mode; and union/intersection/clause structure is
  positive. The classic contravariance counterexample is deflected by
  signature opacity, exactly as the main paper notes for \textbf{L-MONO}.
\item[\textbf{(M2)}] \emph{Union monotonicity.} \textbf{L-UNION} at the
  point of use: $\sem{t_2}_\rho \subseteq \sem{t_2 \vee t}_\rho$.
\end{itemize}
\end{definition}

\begin{lemma}[All Transitions Are Monotone]
\label{lem:trans-mono}
Under M2, every operation of Definition~\ref{def:ext-transitions} except
\texttt{restore} takes $\rho$ to some $\rho'$ with $\rho \stle \rho'$.
\end{lemma}
\begin{proof}
Case analysis on the main paper's transition table plus
Definition~\ref{def:ext-transitions}. Fresh bind
($\mathrm{unbound} \to \SVOpen(t)$ or $\SVClosed(t)$): envelope grows from
$\emptyset$. Fresh covariant get: $\mathrm{unbound} \to \SVClosed(\top)$:
grows to $\sem{\top}$. Open union
($\SVOpen(t_2) \to \SVOpen(t_2 \vee t)$, or
$\SVClosed(t_2 \vee t)$ for non-candidates): grows by M2. Seal
($\SVOpen(t) \to \SVClosed(t)$): envelope unchanged. No-op rows,
\texttt{gen} (the new identifier is unbound before and after — its id was
never allocated before, so no term mentions it), cache inserts,
\texttt{setOrigin}, freeze toggles: envelopes unchanged. No transition
removes or replaces bound material: there is no row that shrinks a bound,
and the main paper's never-widen lemma additionally freezes
$\SVClosed$ rows. Envelopes of \emph{other} variables are untouched by
each write ($\mathsf{values}$ is updated pointwise), though their
envelopes may grow transitively via M1 when their bound mentions
$\alpha$ — still monotone.
\end{proof}

\begin{lemma}[Restore Safety]
\label{lem:restore-safe}
Call a computation segment \emph{well-bracketed} if every
\texttt{restore(s)} it executes uses a snapshot \texttt{s} taken by the
same activation (verified for all fourteen sites in
\S\ref{sec:rollback}). Let a well-bracketed segment start at $\rho_0$ and
execute monotone transitions and restores. Then at every moment,
$\rho_0 \stle \rho_{\mathrm{now}}$.
\end{lemma}
\begin{proof}
Induction over the event sequence. Monotone steps preserve the invariant
by transitivity. A restore jumps to a snapshot taken earlier \emph{within
the segment} (well-bracketing), and the invariant held when the snapshot
was taken, so it holds after the jump.
\end{proof}

\begin{lemma}[Record-Point Coherence]
\label{lem:version-values}
Let a moment $m_0$ (the \emph{record point}) occur at the \emph{entry} of
an activation $A$, with $\rho.\mathsf{version} = v$ at $m_0$, and
consider any later moment $m$ inside $A$'s dynamic extent, under
well-bracketing. Then (i) $\rho.\mathsf{version} \geq v$ at $m$, and (ii)
if $\rho.\mathsf{version} = v$ at $m$, then
$\rho.\mathsf{values} = \rho.\mathsf{values}\text{ at } m_0$.
\end{lemma}
\begin{proof}
(i) The version increases at writes and \texttt{gen}, and decreases only
at restores. Under well-bracketing, every restore executed inside $A$'s
extent uses a snapshot saved inside $A$'s extent, hence saved \emph{after}
$m_0$ ($m_0$ is $A$'s entry), hence with version $\geq v$. So no reachable
restore target has version $< v$.
(ii) By induction on the events between $m_0$ and $m$. All version-$v$
snapshots taken in the extent are taken at moments where, by the
induction hypothesis, \texttt{values} equals its $m_0$ state (a change
since $m_0$ would have bumped the version above $v$, and by (i) it cannot
have come back down except by restoring a version-$v$ snapshot, itself
coherent by induction). A moment with version $v$ is either reached
with no intervening \texttt{values} write since the last coherent
version-$v$ moment, or by restoring a coherent version-$v$ snapshot.
\end{proof}

\begin{finding}[Failed proof attempt, kept for honesty]
\label{find:version-collision}
The naive claim ``\texttt{version} uniquely identifies \texttt{values}
across a whole checking run'' is \emph{false}: restore rewinds the
counter, so divergent exploration branches can reuse a version number
with different bindings (save at $v$, grow to $v{+}1$, restore, grow
differently to $v{+}1$). Any argument that treats the version as a global
state hash — including a reading of the \#446 fix as one — is unsound.
The actual safety of \texttt{tenv.seenSub} comes from the
\emph{combination} of version-keying with path-scoping,
Lemma~\ref{lem:tenv-inert} below.
\end{finding}

%==========================================================================
\section{L-MONO: Refinement-Time State Monotonicity}
\label{sec:lmono}

\textbf{L-MONO} (the main paper's axiom of that name) demands: for every
evolution
$\rho \to \rho'$ performed while refining a condition, and every
$u \in \mathrm{range}(\Env)$,
$v \in \sem{u}_\rho \implies v \in \sem{u}_{\rho'}$.

\subsection{Enumeration of refinement-triggered evolutions}
\label{sec:lmono-enum}

By a walk of every arm of \texttt{evalCond}
(\texttt{TypeCheckRefine.sk:371--541}) and its handlers, the evolutions of
the \emph{main} substitution that condition refinement can trigger are
exactly:

\begin{itemize}
\item[\textbf{(E1)}] \emph{Covariant expansion in the refinement
  primitive.} \texttt{setLocalFilterMap} runs
  \texttt{mapSTypes(\ldots, Some(Covariant()), f)} on the main
  substitution (\texttt{TypeCheckRefine.sk:615}); expansion applies
  \texttt{subst.get(tenv, tvar, Covariant())} to every spine variable
  (\texttt{TType.sk:660}): fresh variables seal to $\top$, open variables
  seal at their bound. Used by \emph{every} rule that factors through the
  primitive.
\item[\textbf{(E2)}] \emph{Union/intersection during residue and clause
  rebuilding.} \texttt{mapSTypes} reassembles clauses with
  \texttt{Clauses::unionN} (\texttt{TType.sk:411}); the
  \texttt{typeof\,==\,"object"} default builds
  \texttt{STypes::union} on the main substitution
  (\texttt{TypeCheckRefine.sk:273--285}). Union itself performs no
  \texttt{get}/\texttt{grow} (verified: no such call in
  \texttt{TTypeUnion.sk}) but recursively calls \texttt{TType::inter} on
  field set-types (\texttt{TTypeUnion.sk:102,124}), reaching E3.
\item[\textbf{(E3)}] \emph{Main-substitution intersections.}
  \texttt{evalIsArray}'s true branch and \texttt{evalInstanceof} run
  \texttt{TObject::inter} on the main substitution
  (\texttt{TypeCheckRefine.sk:334,361}). Inside
  \texttt{TType::inter\_}: covariant expansion of both operands
  (\texttt{TTypeInter.sk:158--159}, seals as in E1); constraint
  attachment on bare-variable operands
  (\texttt{TTypeInter.sk:145--151}, a value operation by
  Lemma~\ref{lem:setconstraint}); speculative concrete
  \texttt{isSub} probes (\texttt{:163--188}); and
  \texttt{TFun::inter}'s two \texttt{TFun::isSub} attempts on the main
  substitution (\texttt{TTypeInter.sk:602, 606}) — on success their
  \texttt{growTVars} growth \emph{commits}, on failure it is restored.
\item[\textbf{(E4)}] \emph{Fresh variable generation.}
  \texttt{evalInstanceof} instantiates the class with fresh inference
  arguments (\texttt{makeTArgs}, \texttt{TEnv.sk:308--323}:
  \texttt{gen} + occurrence-level \texttt{setConstraint} +
  \texttt{setOrigin}); \texttt{newObject} may create synthetic index
  variables (\texttt{TType.sk:1764--1780}).
\item[\textbf{(E5)}] \emph{Cache inserts} in all of the above.
\end{itemize}

Everything else refinement does runs on a \emph{clone}
(\texttt{TypeCheckRefine.sk:89,99,177,197,573}) and does not evolve the
main $\rho$ at all; those sites are \S\ref{sec:lclone}'s subject.

\subsection{Discharge}

\begin{theorem}[L-MONO]
\label{thm:lmono}
Under M1 and M2 (Definition~\ref{def:conditions}), every evolution E1--E5
satisfies the \textbf{L-MONO} conclusion — in fact for \emph{all} type
terms $u$, not only $\mathrm{range}(\Env)$.
\end{theorem}
\begin{proof}
E1--E5 are compositions of the transitions of
Definition~\ref{def:ext-transitions} and internal save/restore pairs
(inside the speculative \texttt{isSub} calls of E3). All fourteen
save/restore sites are well-bracketed (\S\ref{sec:rollback}), so
Lemma~\ref{lem:restore-safe} applies to the whole composition:
$\rho \stle \rho'$. M1 turns this into
$\sem{u}_\rho \subseteq \sem{u}_{\rho'}$.
Two of the main paper's worries dissolve: (a) the
\texttt{setConstraint} threat is void by Lemma~\ref{lem:setconstraint} —
no shared envelope shrinks, so the restriction to
$\mathrm{range}(\Env)$ is not needed; (b) the E3 commits from a
\emph{successful} speculative \texttt{TFun::isSub} are ordinary monotone
grows, so ``exploration that succeeds and keeps its growth'' is safe,
not merely ``exploration that fails and rolls back.''
\end{proof}

\begin{finding}[The sharpest sufficient condition is not needed]
The task of bounding \emph{which} variables refinement may touch
(``refinement only constrains variables introduced by the very expression
being refined'') was pursued and abandoned: it is both false —
\texttt{evalIsArray}/\texttt{evalInstanceof} intersect against
$\Env(x)$'s own object branches, sealing variables shared with the whole
definition's checking — and unnecessary, since sealing and growing are
monotone for \emph{every} sharer. The variable-sharing pattern is
irrelevant once Lemma~\ref{lem:setconstraint} removes the only
non-monotone candidate operation.
\end{finding}

\paragraph{Residual dependence.} Theorem~\ref{thm:lmono} is a genuine
discharge of \textbf{L-MONO} \emph{given the main paper's semantics}: M1
is a property of that semantics (provable by structural induction,
sketched in Definition~\ref{def:conditions}; it deserves a short formal
proof when the semantics section is next revised), and M2 is
\textbf{L-UNION}, already a ledger row. The ledger edit in
\S\ref{sec:status} states this precisely.

%==========================================================================
\section{Monotone Reads Under Growth}
\label{sec:monoreads}

The obligation: a judgement validated under $\rho$ must still hold after
inference evolves $\rho \to \rho'$. For $\SVClosed$ bindings this is
never-widen(iii); the open question is $\SVOpen$ bindings, read without
sealing by $\mathsf{get}(\mathsf{contra})$.

\subsection{The invariant}

\begin{definition}[Read Disciplines]
\label{def:read-disciplines}
\begin{itemize}
\item[\textbf{(MR)}] \emph{Monotone-read invariant.} Every
  contravariant (non-sealing) read is a \emph{supertype-side} read: its
  result is used only as the right-hand side of a validation
  $\sem{\mathrm{lhs}} \subseteq \sem{\cdot}$, so that later growth of the
  variable only enlarges the right-hand side.
\item[\textbf{(BR)}] \emph{Bake-requires-seal.} Every covariant read
  whose result is embedded in a persistent artifact (a memoized result, a
  type installed in an environment) seals the variable, so the baked
  snapshot equals the envelope at all later states
  (never-widen(iii)).
\end{itemize}
\end{definition}

\begin{theorem}[Validated Judgements Are Stable]
\label{thm:stable}
Assume M1, MR, BR, and \textbf{L-SUB}. If \texttt{Clauses::isSub}
accepted $\sem{u_1}_\rho \subseteq \sem{u_2}_\rho$ and
$\rho \to \rho'$ by the transitions of
Definition~\ref{def:ext-transitions} (no restore crossing the
acceptance), then $\sem{u_1}_{\rho'} \subseteq \sem{u_2}_{\rho'}$.
\end{theorem}
\begin{proof}
Decompose the occurrences of variables in $u_1, u_2$ by how validation
read them. Left-side occurrences were read covariantly and are sealed on
the success path (BR); never-widen(iii) freezes their envelopes, so
$\sem{u_1}$ is unchanged... more precisely non-increasing only through
frozen envelopes, hence unchanged. Right-side occurrences were read
contravariantly (MR); their envelopes only grow
(Lemma~\ref{lem:trans-mono}), and by M1 $\sem{u_2}$ only grows. The
accepted inclusion is therefore preserved.
\end{proof}

\subsection{Site audit}
\label{sec:mr-audit}

Every \texttt{expandToClauses}/\texttt{expand}/\texttt{get} site with its
variance, in the pinned revision. Expansion applies its variance uniformly
to every \emph{spine} variable, including through \texttt{TSInst} and
\texttt{TInter} unfolding (\texttt{TType.sk:634--678}); it does not enter
fields or signatures, so no variance flip arises inside a single
expansion — flips happen only where the recursion of \texttt{isSub}
swaps sides.

\begin{center}
\renewcommand{\arraystretch}{1.2}
\small
\begin{tabular}{@{}>{\raggedright\arraybackslash}p{5.4cm}p{1.1cm}>{\raggedright\arraybackslash}p{6.7cm}@{}}
\toprule
site & var. & discipline check \\
\midrule
\texttt{isSub\_} lhs $u_1$ (TTypeSubtyping.sk:1623) &
  cov & BR: seals; memo \texttt{seenSub} baked on success \\
\texttt{isSub\_} rhs $u_2$ (:1624) &
  contra & MR: supertype side \\
\texttt{TFun::isSub\_} rest, \emph{sub} side (:1708) &
  contra & MR: becomes rhs of param check
  \texttt{isSub(p2,\,p1)} \\
\texttt{TFun::isSub\_} rest, \emph{super} side (:1768--1773) &
  cov & BR: lhs of param checks; sealed \\
\texttt{TTypeInter} both operands (TTypeInter.sk:158--159) &
  cov & BR: result memoized/installed; seal justifies
  \texttt{seenInter} \\
\texttt{TypeCheckUtils} truthiness reads (:19,43) &
  cov & BR: result drives diagnostics/refinement \\
\texttt{resolveSpreadToTuple} (TVariadicTyping.sk:53,74) &
  contra & MR: tuple shape probe; restored when \texttt{None}
  (:64,70,78) \\
TEnv.sk:820,2103; TypeCheckClass.sk:55,146;
  TypeCheckExpr.sk:1549--1553 & cov & BR \\
\texttt{expandTVarsForPp} (TType.sk:703) & cov &
  on a \emph{clone} (TTypePp.sk:815, TTypeSubtyping.sk:1195--1196):
  pretty-printing must not seal \\
\texttt{ArrayLitTyping} probes (TArrayLitTyping.sk:175,236,394) &
  cov & :236 clones; :175,394 seal on the live subst
  (commitment sites) \\
\bottomrule
\end{tabular}
\end{center}

No site violates MR: all contravariant reads end up on the right of a
validation. BR holds at all sites on the success path, with one
exception, which is the section's finding.

\subsection{Residual hazard: the freeze-window unsealed read}
\label{sec:freeze-hazard}

$\mathsf{get}(\mathsf{cov})$ on an open \emph{candidate} under
\texttt{freezeCandidates} returns the bound \emph{without sealing}
(\texttt{TTypeSubst.sk:202--203}) — a covariant read that violates BR by
design, so that constraint admissibility checks are queries. The check
runs full \texttt{TType::isSub(boundTy, cstr)}
(\texttt{TTypeSubtyping.sk:1435}), whose \emph{successes} insert pairs
into the persistent, version-\emph{unkeyed} \texttt{subst.seenSub} memo.

\begin{finding}[Stale \texttt{seenSub} memo via freeze-window reads]
\label{find:freeze-memo}
Candidate scenario: a variable $\alpha$ (an ordinary candidate,
$\SVOpen(t)$) occurs in \emph{both} sides of a check — e.g.\
$u_1 = \mathrm{Box}\langle\alpha\rangle$ and $u_2$ mentions a constrained
variable whose constraint check validates
$\mathrm{isSub}(u_1, \mathrm{cstr})$ under freeze, reading $\alpha$ at
$t$ \emph{unsealed} and memoizing $(u_1, \mathrm{cstr}) \in
\mathsf{seenSub}$. After the freeze pops, \texttt{growUncovered} may
contravariantly grow the same $\alpha$ (it is no longer protected) to
$t \vee s$. The memo now asserts an inclusion validated with the smaller
$\sem{u_1}$; a later syntactically-identical query hits it
(\texttt{TTypeSubtyping.sk:1588}) without re-validation. The eventual
covariant expansion of $u_1$ seals $\alpha$ at $t \vee s$ — \emph{after}
the memo was recorded. The discipline itself cannot close this gap: the
memo survives on the enclosing check's success path, and nothing re-keys
it by version. Mitigations narrow the window without closing it (memos
recorded under a speculative subtree are erased by its restore, since
$\mathsf{seenSub}$ is in \texttt{SubstState}; the scenario needs the
same open candidate reachable from both sides plus an exact term-level
re-query). The gap is real: an end-to-end witness exists — an
order-dependent unsound acceptance in which \emph{adding} an earlier call
to a generic function (validating its constraint under freeze against the
still-open candidate, and memoizing the success) flips a correct
contravariance rejection of the post-widening call into an acceptance,
so a \texttt{number}-typed binding holds \texttt{"s"} at runtime.
MR-1 (\S\ref{sec:status}) is therefore \textsc{violated}. Fix: key
\texttt{subst.seenSub} by version, or suppress memo insertion while
\texttt{freezeCandidates} is set; the same fix should cover
\texttt{seenInter}, which the witness also exercises. Issue draft:
\path{issues/freeze-window-stale-memo-caches.md}; a pinned test is
recommended.
\end{finding}

Theorem~\ref{thm:stable} discharges the obligation
\emph{conditional on} MR-1 and on \textbf{L-SUB}, and the audit above
discharges its MR/BR premises for every other site; since MR-1 fails, the
obligation stands \textsc{violated} exactly at the freeze-window memo and
\textsc{proved} everywhere else.

%==========================================================================
\section{L-CLONE: Clone--Compute--Install}
\label{sec:lclone}

\subsection{The pattern and what is actually lost}

\texttt{evalVarEquality} computes
$u = \mathsf{inter}(\Env(x), \tau_e)$ under \texttt{subst.clone()} and
installs $u$ into the scope, discarding the clone
(\texttt{TypeCheckRefine.sk:87--94}); likewise \texttt{removeConstant}
(\texttt{:96--106}), \texttt{evalFieldEqualityType}
(\texttt{:176--201}), and \texttt{evalFieldTruthiness}
(\texttt{:571--576}). \texttt{clone()} forks all seven
\texttt{SubstState} fields (immutable maps, so the fork is a true
snapshot) and shares only \texttt{results} and \texttt{context}
(\texttt{TTypeSubst.sk:137--149}).

Write $\rho_f$ for the fork state ($=$ the main state at the clone),
$\rho_c$ for the clone state when \texttt{inter} returns, and $\rho''$
for any later main state. What the clone loses at discard, classified by
what \texttt{inter} can do:

\begin{itemize}
\item[\textbf{(C1)}] \emph{Constraint attachments}
  (\texttt{TTypeInter.sk:145--151}). By
  Lemma~\ref{lem:setconstraint} these never touched the clone's state at
  all: the constraint lives in the returned occurrence inside $u$, which
  is installed. \emph{Nothing is lost.} A natural conjecture
  (``if constraints live in the returned type, clone-discard may be
  benign for constraints but not for value bindings'') is confirmed on
  its first half.
\item[\textbf{(C2)}] \emph{Seals} from covariant expansion
  (\texttt{TTypeInter.sk:158--159}): the expansion replaces the spine
  occurrence by the read value in $u$ — fresh variables bake $\top$,
  open variables bake their fork-state bound $t$ — but the seal itself
  dies with the clone; on the main substitution the variable stays
  open/unbound and may later grow or seal differently.
\item[\textbf{(C3)}] \emph{Grows and commitments} from speculative
  subtyping inside \texttt{inter}: \texttt{TFun::inter} runs
  \texttt{TFun::isSub} both ways on the clone
  (\texttt{TTypeInter.sk:602, 606}), committing growth on success;
  \texttt{ArrayLitTyping} may commit a literal's form variable on the
  clone. All lost at discard, while $u$ may mention the affected
  variables.
\item[\textbf{(C4)}] \emph{Fresh generations} on the clone (expansion of
  a cold \texttt{TSInst} can reach \texttt{tenv.instantiate} and, through
  class setup, \texttt{gen}-sites such as synthetic index fields,
  \texttt{TType.sk:1764--1780}). The clone's version counter advances
  privately; the main counter does not.
\end{itemize}

\subsection{Over-approximation: the positive result}

\begin{theorem}[Clone Stability]
\label{thm:lclone}
Assume M1, M2, and the fork-state form of \textbf{L-INTER}:
\begin{itemize}
\item[\textbf{(I$_f$)}] $\sem{t_1}_{\rho_f} \cap \sem{t_2}_{\rho_f}
  \subseteq \sem{u}_{\rho''}$ \emph{for every $\rho'' \sqsupseteq
  \rho_f$}, where $u = \mathsf{inter}$ computed by a well-bracketed run
  from $\rho_f$.
\end{itemize}
Then \textbf{L-CLONE} holds: the installed $u$ over-approximates the
fork-state intersection under the main substitution at install time and
at all later states. Moreover I$_f$ follows from the ordinary
\textbf{L-INTER} plus M1 for every value whose membership does not
depend on clone-side state that $u$ bakes or mentions; the two ways such
dependence can arise are exactly C2/C3 occurrences and C4 identifiers,
analyzed below.
\end{theorem}
\begin{proof}[Proof sketch]
The clone evolves monotonically from $\rho_f$
(Lemma~\ref{lem:trans-mono} + Lemma~\ref{lem:restore-safe}):
$\rho_f \stle \rho_c$. The main substitution also evolves monotonically:
$\rho_f \stle \rho''$. Ordinary \textbf{L-INTER} gives containment in
$\sem{u}_{\rho_c}$; the gap is that $\rho_c$ and $\rho''$ are
\emph{incomparable} (divergent branches of $\rho_f$), so M1 cannot bridge
them directly — this was the first failed proof attempt, and the reason
\textbf{L-CLONE} cannot be reduced to \textbf{L-INTER} + M1 alone.
The bridge is per-occurrence: (C2) material \emph{baked} into $u$ is the
clone-state envelope of the variable at read time, which by monotonicity
of the clone run is $\supseteq$ its fork-state envelope; baked material
is state-independent thereafter, so reading $u$ under $\rho''$ yields at
least what reading it with fork-state envelopes would — and the
fork-state intersection is contained in that by \textbf{L-INTER}
\emph{evaluated at} $\rho_f$. (C1) occurrences carry their constraints
with them and denote
$\env{\rho''}{\alpha} \cap \bigcap\sem{\mathrm{cstr}}$, which by
$\rho_f \stle \rho''$ contains the fork-state value — the intersection
semantics is exact here. (C3) growth affects $u$ only through signature
positions (\texttt{TFun::inter} returns one of the two function terms,
whose parameter/return types may mention clone-grown variables), and
signatures are opaque to $\sem{\cdot}$ — the main paper's
signature-opacity deflection applies verbatim. Field-get positions reach
\texttt{inter} recursively and fall under C1/C2. What the sketch does
\emph{not} cover is C4, next.
\end{proof}

\subsection{Clone-escaped fresh identifiers: an unreachable hazard}

\begin{finding}[Identifier collision between clone and main is
unreachable]
\label{find:clone-gen}
Fresh identifiers are drawn from the \emph{version counter}
(\texttt{TTypeSubst.sk:237--253}), and a clone forks the counter: if a
clone ever generated a variable (C4), its id would equal the id the
\emph{main} substitution assigns to its own next generation, and a
clone-generated occurrence escaping into the installed $u$ in a
semantically visible position (a field get type — synthetic index fields
qualify, \texttt{TType.sk:1764--1780}) would, once the main substitution
binds the colliding id, silently read an unrelated variable's envelope.
Nothing in Theorem~\ref{thm:lclone}'s bridge would cover that: the
occurrence's denotation would not be monotone from fork state. The hazard
is settled by a reachability audit: of all ten \texttt{gen*} call sites,
\emph{none} lies in the clone-compute-install call tree
(\texttt{inter} / \texttt{removeConstant} / \texttt{expand} /
\texttt{instantiate} / \texttt{normalize} / \texttt{newObject} /
\texttt{addFields}); in particular the two
\texttt{createSyntheticIndex} generators have exactly two callers, both
in the object-literal branch of expression typing, which runs on the
\emph{main} substitution — never under a refinement clone. So no
clone-side fresh identifier is ever generated during
clone-compute-install, no collision can occur, and CL-1
(\S\ref{sec:status}) is \emph{discharged}. The discharge rests on a
code invariant that should be guarded: assert (and CI-check) that
clone-side runs of \texttt{inter} never call \texttt{gen*}; drawing ids
from a counter shared by reference would remove the hazard class
altogether.
\end{finding}

\paragraph{Shared \texttt{results}.} Clone-side speculative
\texttt{TFun::isSub} successes record subtyping edges into the shared
accumulator even when the exploration is discarded
(\texttt{TTypeSubtyping.sk:1681--1688}). This does not touch typing
state; it can only over-connect the callgraph consumed by later phases.
Out of scope for type soundness, recorded here so the boundary is
explicit.

%==========================================================================
\section{No Commitment on Failure: the Rollback Audit}
\label{sec:rollback}

The governing contract is stated at the top of the subtyping module
(\texttt{TTypeSubtyping.sk:4--7}): \emph{every \texttt{isSub} function
must guarantee that when it returns \texttt{STError}, the subst is
unchanged}. This section verifies it, pairs every save with its restores,
and documents the one place the letter of the contract is violated by
design.

\subsection{Exhaustive save/restore pairing}

All \texttt{save}/\texttt{restore} sites in the pinned revision
(\texttt{grep} over \texttt{src/}), each verified well-bracketed
(restore uses the same activation's snapshot; no early exit escapes the
pair):

\begin{center}
\renewcommand{\arraystretch}{1.15}
\small
\begin{tabular}{@{}p{0.5cm}>{\raggedright\arraybackslash}p{7.4cm}>{\raggedright\arraybackslash}p{5.4cm}@{}}
\toprule
\# & site (save; restore) & restore condition \\
\midrule
1 & \texttt{TFields::isSub} (\texttt{TTypeSubtyping.sk:796; 798}) &
  on \texttt{STError} \\
2 & \texttt{STypes::isSub} (\texttt{:1109; 1136}) & on failure \\
3 & \texttt{Clauses::isSub} (\texttt{:1149; 1169}) &
  on overall failure (clause with no matching partner) \\
4 & \texttt{TObject::isSub} outer (\texttt{:1877; 1889}) &
  on overall failure \\
5 & \texttt{TObject::isSub} per branch pair (\texttt{:1881; 1886}) &
  after each failed pairing, before trying the next \\
6 & \texttt{TType::isSub\_} (\texttt{:1617; 1643}) &
  on validation failure (growth persists on success, intentionally) \\
7 & \texttt{growTVars} per-variable (\texttt{:1416; 1447}) &
  on constraint-check failure \\
8 & \texttt{isCovered} probes (\texttt{:1354; 1365}) &
  after \emph{every} probe, hit or miss \\
9 & \texttt{TFun::isSub} (\texttt{:1676; 1679}) & on \texttt{STError} \\
10 & \texttt{ArrayLit} tuple commit (\texttt{TArrayLitTyping.sk:402;
  416}) & on element-check failure (case 2c) \\
11 & \texttt{resolveSpreadToTuple} (\texttt{TVariadicTyping.sk:51;
  64,70,78}) & on every \texttt{None} path \\
12 & \texttt{inferRestTypeVar} (\texttt{TVariadicTyping.sk:124; 130}) &
  on constraint failure (error also reported) \\
13 & \texttt{satisfies} (\texttt{TypeCheckExpr.sk:2793; 2805}) &
  \emph{unconditionally} — \texttt{satisfies} must not commit inference \\
14 & clone-as-rollback: \texttt{TEnv.sk:222},
  \texttt{TypeCheckRefine.sk} (5 sites), \texttt{TypeCheckExpr.sk}
  (6 sites), pretty-printer & discard instead of restore \\
\bottomrule
\end{tabular}
\end{center}

Three structural corollaries. (i) Well-bracketing holds at all sites, so
Lemmas~\ref{lem:restore-safe} and~\ref{lem:version-values} apply
globally. (ii) \texttt{SavedSubst} captures \emph{all} mutable typing
state (Definition~\ref{def:fullstate}); the two exclusions are
\texttt{freezeCandidates} — safe by its own push/pop bracketing, U1 —
and \texttt{tenv.seenSub} — safe by \S\ref{sec:tenv-seensub}. In
particular a rollback also erases \texttt{seenSub}/\texttt{seenInter}/%
\texttt{seenUnion}/\texttt{seenExpansion} entries recorded in the failed
subtree, so no stale \emph{success} memo survives a rollback. (iii)
Because no environment binding is written between any save and its
restore (refinement installs happen strictly after the operations
return), a restore reinstates a state under which every installed type
was already satisfied — rollback needs no analogue of \textbf{L-MONO}.

\subsection{The intentional exception: array-literal commitments}
\label{sec:arraylit-exception}

\texttt{ArrayLitTyping::subtyping} runs at the top of
\texttt{TType::isSub\_} \emph{before} the save at \texttt{:1617}
(\texttt{TTypeSubtyping.sk:1612}). Its cases 1 and 2b commit the
literal's form variable and then return the result of a subsequent
check \emph{even when that check fails}
(\texttt{TArrayLitTyping.sk:334--341, 368--392}): the commitment
persists on an \texttt{STError} return, violating the module contract's
letter. This is deliberate (``irreversible'', \texttt{:140}); case 2c by
contrast restores the commitment on failure (row 10).

\begin{lemma}[The carve-out is sound for accepted programs]
\label{lem:arraylit-carveout}
A commitment persisting on a failed check cannot cause unsound
acceptance, provided every caller that \emph{discards} an
\texttt{STError} (uses \texttt{isSub} as a probe) either wraps the call
in save/restore or cannot reach an uncommitted array literal.
\end{lemma}
\begin{proof}[Proof, with caller audit]
On a non-discarded failure the program is rejected — the committed type
may be wrong for the runtime value (case 2b commits to
\texttt{Array<expected>} even when the elements fail against
\texttt{expected}), but no judgement about an accepted program uses it.
The risk is a discarded failure. Probe callers: \texttt{isCovered}
(row 8) restores around every probe; \texttt{growTVars} constraint checks
(row 7) restore on failure; the speculative probes inside
\texttt{inter\_} (\texttt{TTypeInter.sk:163--188}) and
\texttt{SType::overlaps} (\texttt{:392--393}) compare a boxed-primitive
instance against a single \texttt{SType} — neither side is a bare
uncommitted \texttt{TArrayLit}, which is the only trigger shape
(\texttt{getSingleArrayLitUncommitted} requires a sole
\texttt{TArrayLit} inst, \texttt{TArrayLitTyping.sk:68--82});
\texttt{TFun::inter}'s two attempts run under \texttt{TFun::isSub}'s own
save/restore (row 9), which was taken \emph{before} the inner
\texttt{isSub\_} ran, hence before the commitment — a failed first
attempt rolls the commitment back. The constraint probe at
\texttt{TEnv.sk:222} runs on a clone (row 14).
\end{proof}

This audit is complete for the pinned revision but fragile — any new
\texttt{isSub}-as-probe without a save/restore wrapper re-opens it.
Sub-obligation RB-1 in \S\ref{sec:status} records the maintenance
condition.

\subsection{\texttt{tenv.seenSub}: the \#446 argument, made precise}
\label{sec:tenv-seensub}

\texttt{tenv.seenSub} is the coinduction set of \texttt{isSub\_}: on
entry, if $(\mathsf{version}, u_1, u_2)$ is present the check succeeds by
coinductive hypothesis; otherwise the triple is recorded and threaded
into the recursive calls (\texttt{TTypeSubtyping.sk:1603--1604}). It is
not part of \texttt{SubstState}, and no rollback ever touches it. Its
safety rests on two independent mechanisms:

\begin{lemma}[Version-Keyed Hypotheses Are Inert or Exact]
\label{lem:tenv-inert}
A lookup that hits an entry $(v, u_1, u_2)$ observes
$\mathsf{values}$ identical to the recorder's state at record time.
Consequently a hypothesis is either consulted under exactly the binding
state it was made in, or (after any growth, seal, or fresh generation,
each of which bumps the version) it is \emph{inert} — present but
unmatchable.
\end{lemma}
\begin{proof}
Two mechanisms. \emph{Path-scoping} (U2, \S\ref{sec:model}):
\texttt{TEnv} is frozen, so the entry is visible only in the dynamic
extent of the recording activation — entries never leak across sibling
explorations or survive returns, which already rules out the
cross-branch collision of Finding~\ref{find:version-collision}.
\emph{Record-point coherence}: the record happens at the activation's
entry, so Lemma~\ref{lem:version-values} applies to every lookup in the
extent: version $= v$ implies \texttt{values} equal to record time.
(The lemma's premise — restores in the extent use snapshots from the
extent — is row-by-row verified in \S\ref{sec:rollback}.)
\end{proof}

This is the precise content of the \#446 regression class and its guard:
without the version key, a rolled-back exploration's hypothesis could be
consulted under different bindings \emph{within the same extent} (probe
loops, rows 7--8, restore and re-grow inside one \texttt{isSub\_}
activation); the version key makes exactly those entries inert. The
prerequisite ``version is bumped on every \texttt{values} write'' is
verified at all three write sites (\texttt{TTypeSubst.sk:181,195,205});
note \texttt{gen*} bumps too, which additionally protects hypotheses
against fresh-variable aliasing inside an extent.

By contrast, \texttt{subst.seenSub} — the success memo, \emph{not}
version-keyed — is protected against rollback staleness by living inside
\texttt{SubstState} (corollary (ii) above), and against growth staleness
by Theorem~\ref{thm:stable}; it is the memo that inherits the MR-1
hazard (Finding~\ref{find:freeze-memo}).

%==========================================================================
\section{Status Table and Ledger Edits}
\label{sec:status}

Named sub-obligations introduced by this paper:

\begin{itemize}
\item[\textbf{MR-1}] No \texttt{subst.seenSub} memo recorded while
  \texttt{freezeCandidates} is set may outlive subsequent growth of an
  open candidate readable from its left-hand side
  (Finding~\ref{find:freeze-memo}). Status: \textsc{violated}, witnessed
  end-to-end (issue draft
  \path{issues/freeze-window-stale-memo-caches.md}). Fix options:
  version-key the memo, or skip memo insertion under freeze; either must
  also cover \texttt{seenInter}. Pinned test recommended.
\item[\textbf{CL-1}] No fresh variable generated on a clone may occur in
  a type installed against the main substitution
  (Finding~\ref{find:clone-gen}). Status: discharged — the property
  holds, verified unreachable by the \texttt{gen*} call-site audit;
  guard the underlying invariant with the recommended CI assertion
  (\texttt{gen*} unreachable under clones used by refinement).
\item[\textbf{RB-1}] Maintenance condition: every \texttt{isSub} caller
  that discards \texttt{STError} must wrap the call in save/restore or
  be unreachable by uncommitted array literals
  (Lemma~\ref{lem:arraylit-carveout}); a CI grep for new probe sites is
  the cheap enforcement.
\item[\textbf{M1}] Semantic monotonicity of $\sem{\cdot}$ in variable
  envelopes (Definition~\ref{def:conditions}) — a property of the main
  paper's semantics, provable by structural induction; should be added
  there as a lemma rather than assumed here.
\end{itemize}

What this paper establishes, row by row:

\begingroup
\renewcommand{\arraystretch}{1.25}
\small
\begin{tabular}{@{}>{\raggedright\arraybackslash}p{4.2cm}>{\raggedright\arraybackslash}p{9.4cm}@{}}
\toprule
\textbf{Row} & \textbf{Established here} \\
\midrule
L-MONO & Theorem~\ref{thm:lmono}, conditional on M1
  and L-UNION; the \texttt{setConstraint} threat is void
  (Lemma~\ref{lem:setconstraint}) and the $\mathrm{range}(\Env)$
  restriction is droppable \\
Obl.\ Monotone Reads & general statement via
  Theorem~\ref{thm:stable} (given M1 + the MR/BR audit,
  \S\ref{sec:mr-audit}); the residue MR-1 is witnessed
  (Finding~\ref{find:freeze-memo}) \\
L-CLONE & Theorem~\ref{thm:lclone} given I$_f$
  (fork-state L-INTER; registry row I-F); the id-collision hazard CL-1
  discharged as
  unreachable (Finding~\ref{find:clone-gen}), CI-assertion guard
  recommended \\
Obl.\ No Commitment on Failure & 14-row save/restore
  audit (\S\ref{sec:rollback}; \texttt{SavedSubst} completeness;
  \texttt{freezeCandidates} and \texttt{tenv.seenSub} exclusions
  justified, Lemma~\ref{lem:tenv-inert}); array-literal commitments
  persist on failure by design, sound by
  Lemma~\ref{lem:arraylit-carveout} under RB-1 \\
\bottomrule
\end{tabular}
\endgroup

\paragraph{Registry.}
The authoritative, current statuses are maintained in the registry
(\texttt{docs/ledger.yaml}), rendered as the main paper's Assumptions and
Proof Obligations tables; this paper is cited there as evidence. The
registry rows it supports:

% GENERATED by docs/tools/docs.py from docs/ledger.yaml — do not edit.

\begin{center}
\renewcommand{\arraystretch}{1.25}
\small
\begin{tabular}{@{}>{\raggedright\arraybackslash}p{2.4cm}>{\raggedright\arraybackslash}p{8.2cm}>{\raggedright\arraybackslash}p{2.6cm}@{}}
\toprule
\textbf{Row} & \textbf{Statement} & \textbf{Status} \\
\midrule
L-MONO & Refinement-time state evolutions preserve membership in already-installed types. & \textsc{proved} \\
L-CLONE & A type computed by \texttt{inter} under a cloned substitution over-approximates the intersection under the main one. & \textsc{proved} \\
OBL-MONOTONE & Monotone Reads Under Growth: judgements validated under $\rho$ survive later legal transitions. & \textsc{violated} \\
OBL-RESTORE & No Commitment on Failure: failing and exploratory paths restore all relevant state. & \textsc{proved} \\
MR-1 & No \texttt{subst.seenSub} memo recorded under \texttt{freezeCandidates} may outlive growth of an open candidate readable from its left side. & \textsc{violated} \\
CL-1 & No fresh variable generated on a clone occurs in a type installed against the main substitution. & \textsc{proved} \\
M1 & Semantic monotonicity of $\sem{\cdot}$ in variable envelopes (a property of the get-only, signature-opaque semantics). & \textsc{assumed} \\
I-F & Fork-state form of L-INTER: the containments hold from a forked substitution state (consumed by L-CLONE's proof). & \textsc{assumed} \\
RB-1 & Maintenance condition: \texttt{isSub} callers discarding errors are save/restore-wrapped or unreachable by uncommitted array literals. & \textsc{open} \\
SETCONSTRAINT-FUN & \texttt{TVar::setConstraint} is a functional update on a frozen value: constraints are occurrence-level, not shared state. & \textsc{proved} \\
E-TRANSPORT & Monotone-reads transport: an inclusion validated under $\rho$ survives later legal transitions on the same timeline, modulo MR-1. & \textsc{proved} \\
TENV-EXPANDINTER & The \texttt{tenv.expandInter} mode flag is path-scoped, and both intersection forms it selects (expanded and deferred \texttt{unexpandedInter}) satisfy the L-INTER containment. & \textsc{open} \\
\bottomrule
\end{tabular}
\end{center}


\paragraph{Imports (review contract).}
The named statements this paper assumes from elsewhere. A reviewer treats
each as an axiom, checking citations against the registry's canonical
statement; no other paper's internals are in scope:

% GENERATED by docs/tools/docs.py from docs/ledger.yaml — do not edit.

\begin{center}
\renewcommand{\arraystretch}{1.25}
\small
\begin{tabular}{@{}>{\raggedright\arraybackslash}p{2.6cm}>{\raggedright\arraybackslash}p{7.4cm}>{\raggedright\arraybackslash}p{1.2cm}>{\raggedright\arraybackslash}p{2.2cm}@{}}
\toprule
\textbf{Row} & \textbf{Canonical statement} & \textbf{Owner} & \textbf{Status} \\
\midrule
L-UNION & \texttt{union} over-approximates set union (all levels). & B & \textsc{proved} \\
L-INTER & \texttt{inter} over-approximates set intersection (the containment clause). & B & \textsc{proved} \\
L-SUB & \texttt{Clauses::isSub} acceptance implies semantic inclusion under the current state. & D & \textsc{open} \\
\bottomrule
\end{tabular}
\end{center}


\paragraph{Summary.} Three of the four rows close (L-MONO; No
Commitment on Failure, with a documented carve-out; L-CLONE modulo
I$_f$, its id-collision hazard verified unreachable and guarded by a CI
assertion). The fourth — Monotone Reads — is \textsc{proved} at every
audited site except its one sharp residue, MR-1, which is
\textsc{violated} with an end-to-end witness (issue draft
\path{issues/freeze-window-stale-memo-caches.md}); the fix
(version-keying or freeze-window write suppression) must cover both
\texttt{seenSub} and \texttt{seenInter}.

%==========================================================================
\appendix
\section{Implementation Anchors}

Pinned to \texttt{\skjspinhash}.

\begin{center}
\renewcommand{\arraystretch}{1.25}
\small
\begin{tabular}{@{}>{\raggedright\arraybackslash}p{4.8cm}@{\quad}>{\raggedright\arraybackslash}p{9.8cm}@{}}
\toprule
\textbf{Model} & \textbf{Implementation} \\
\midrule
$\rho$ (seven fields) & \texttt{SubstState} (TTypeSubst.sk:13--49) \\
snapshot / reinstate & \texttt{SavedSubst}, \texttt{save},
  \texttt{restore} (TTypeSubst.sk:36--74) \\
transitions & \texttt{Subst::grow/get} (TTypeSubst.sk:151--222) \\
fresh ids from version & \texttt{gen/genForField/genForLiteralCommitment}
  (TTypeSubst.sk:237--253) \\
$\mathsf{freeze}$ (unsaved) & \texttt{Subst.freezeCandidates}
  (TTypeSubst.sk:76--89); push/pop TTypeSubtyping.sk:1431--1443 \\
clone fork & \texttt{Subst::clone} (TTypeSubst.sk:137--149) \\
occurrence constraints & \texttt{TVar}, \texttt{setConstraint}
  (TType.sk:2576--2607) \\
success memo & \texttt{subst.seenSub}
  (TTypeSubtyping.sk:1588,1591) \\
coinduction hypotheses & \texttt{tenv.seenSub} (TEnv.sk:74;
  TTypeSubtyping.sk:1603--1604); \texttt{tenv.seenInter}
  (TEnv.sk:69; TTypeInter.sk:154--157) \\
result memos & \texttt{seenInter/seenUnion/seenExpansion}
  (TTypeInter.sk:107--125; TTypeUnion.sk:29--32; TType.sk:650--653) \\
uniform-variance expansion & \texttt{expandToVectors}
  (TType.sk:634--678) \\
validation spine & \texttt{isSub\_/growTVars/isCovered}
  (TTypeSubtyping.sk:1318--1646) \\
clone-compute-install & \texttt{evalVarEquality} etc.\
  (TypeCheckRefine.sk:87--106,176--201,571--576) \\
main-subst refinement ops & \texttt{setLocalFilterMap}
  (TypeCheckRefine.sk:615); \texttt{evalIsArray/evalInstanceof}
  (TypeCheckRefine.sk:334,361) \\
form commitments & \texttt{ArrayLitTyping}
  (TArrayLitTyping.sk:140--148,326--432) \\
external grow callers & TArrayLitTyping.sk:147,352;
  TVariadicTyping.sk:125; TypeCheckUtils.sk:310,321,332 \\
\bottomrule
\end{tabular}
\end{center}

\end{document}
